\section{The unipotent period map}\label{sec:unipotent}

We construct a map from semi-infinite cohomology to the period of
localization.  Its underlying map is the inclusion of constant functions
on a framed bundle space.  The main point is to identify the complexes
on both sides with their correct determinant lines, and to make this map
compatible with changing the marked divisor and its formal lift.
The character-twisted period construction of
\cite[Section~12.5.17 and Theorem~12.5.18]{GLCII}, and its extension to
reductive-group inputs in \cite[Theorem~12.8.8]{GLCII}, provide the
corresponding Ran comparison.  Here we construct the map before Ran
integration, retaining the relative and level compatibilities needed
for the family comparison.

\subsection{The normalization and the statement}

Keep the marked divisor $E$, the chosen level divisor, and the
Harish--Chandra categories of Section~\ref{sec:localization}.  Write
$h_E^0$ for the full $N_0$ arc Lie lattice and
$h=h_E^{\mathrm{lv}}$ for its level sublattice.  On either the $N$ or
the $G$ input category, define
\[
 B_E^0(M)=C^{\infty/2,h_E^0}
       \bigl(\operatorname{Lie}LN_E,M\otimes k_{-\chi}\bigr),
 \qquad
 B_E^{\mathrm{lv}}(M)=C^{\infty/2,h}
       \bigl(\operatorname{Lie}LN_E,M\otimes k_{-\chi}\bigr).
\]
These are the continuous extensions of the compact-object calculations
in \eqref{eq:L6}; additional structures are forgotten.  As in
Section~\ref{sec:loop}, after applying the forgetful functor $b$, we write the underlying
relative formulas over the smooth affine parameter scheme $S$ without
angle brackets.  Their crystalline versions on parameter
schemes use the angle-bracket convention.  The full-arc formula
$B_E^0$ has the latter interpretation directly; the level formula
requires the following relative determinant normalization.

Put
\[
 m=\operatorname{rank}(h_E^0/h),\qquad
 l^{\mathrm{lv}}=\det R\Gamma(X,\mathfrak n^{\mathrm{lv}}),\qquad
 \delta=-\chi(X,\mathfrak n^{\mathrm{lv}}).
\]
The exact sequence for the level divisor, together with change of
semi-infinite lattice, gives
\begin{equation}\label{eq:T9}
 \begin{aligned}
 \delta&=\delta_N+m,\\
 l^{\mathrm{lv}}
   &=l_N\otimes\det(h_E^0/h)^{-1},\\
 B_E^{\mathrm{lv}}
   &=B_E^0\otimes\det(h_E^0/h)[m].
 \end{aligned}
\end{equation}
These identities also hold for intermediate tuples: the quotient
$h_E^0/h$ depends only on the level divisor.

\begin{proposition}[Period map]\label{prop:unipotent-period}
For the $N$ Harish--Chandra category there is a natural continuous map
\begin{equation}\label{eq:T10}
 l^{\mathrm{lv}}[\delta]\otimes B_E^{\mathrm{lv}}
       \longrightarrow \operatorname{Per}_{N,E}^{\mathrm{lv}}.
\end{equation}
After shifting by $[-2\delta]$, its source identifies canonically with
$l_N[-\delta_N]\otimes B_E^0$.  In this normalization the map is a
map of crystalline operations, compatible with the parameter base
changes, lift comparisons, and deeper-level comparisons of
Section~\ref{sec:localization}.
\end{proposition}

We first construct the map on bounded ordinary modules over an affine
base.  Finite-dimensional localization will identify its target with
semi-infinite cohomology having an extra function-space coefficient.
We then construct the same comparison on the compact presentations of
the renormalized categories.  This second construction proves the
derived base-change assertion without an exactness assertion for
invariants under a fixed congruence subgroup.

\subsection{A finite presentation of the bundle space}

Abbreviate the bundle stack $B_N^{\mathrm{lv}}$ to $B_N$, and let
$\widetilde B_N$ denote its space of full frames along $E$.  Framed
replacement gives it a right $N$-loop action.  Specifying the full
frame already specifies the level structure, so $\widetilde B_N$ is
independent of level: it is also the space of ordinary $N_0$-bundles
with frames on these disks.

Use the height lattices $L_l$ from the proof of \eqref{eq:T8}, now
along $E$.  Enlarge $l$ until the first cohomology of every enlarged
root line vanishes, relatively over the affine base.  Set
\[
 H=H_{N,E}^{\mathrm{lv}},\qquad
 A_l=\operatorname{Lie}L_l.
\]
Let $K_l$ be the group of global sections of the enlarged unipotent
group scheme and put $k_l=\operatorname{Lie}K_l$.  It is a smooth
finite-dimensional unipotent group over the base.

\begin{lemma}\label{lem:unipotent-presentation}
Disk gluing identifies
\begin{equation}\label{eq:T11}
 B_N=[K_l\backslash L_l/H],\qquad
 \widetilde B_N=K_l\backslash L_l.
\end{equation}
For $e_l=\operatorname{rank}(A_l/h)$ and
$f_l=\operatorname{rank}k_l$, one has
\begin{equation}\label{eq:T12}
 \delta=e_l-f_l,\qquad
 l^{\mathrm{lv}}=\det(k_l)\otimes\det(A_l/h)^{-1}.
\end{equation}
The character $\chi$ on the presentation induces the function
$\psi_0$ on $B_N$.
\end{lemma}

\begin{proof}
Filter the enlarged group by root height.  Each successive quotient
is a vector group whose underlying vector bundle has vanishing
$H^1$.  Its torsors on $X$ therefore become trivial over an affine
chart of the parameter base.  Induction along this filtration
trivializes every enlarged-group torsor; two trivializations differ
by $K_l$.  Recovering the original bundle amounts to choosing its
lattice at $E$, which gives the double quotient in \eqref{eq:T11}.
With full frames the right quotient by $H$ is omitted.

The inclusion of global sections into sufficiently long jets is
fiberwise injective.  After increasing the jet order on the chosen
base chart, it realizes $K_l$ as a closed smooth subgroup of a finite
quotient of $L_l$.  Exponential coordinates compatible with root
height identify the successive quotient maps with affine-space
torsors.  Thus the quotients occurring here have smooth relative
scheme presentations.  Enlarging $l$ recovers the same framed space
with its right action.

For each root bundle the meromorphic-lattice sequence represents
$R\Gamma(X,\mathfrak n^{\mathrm{lv}})$ by global sections of the
enlarged bundle mapping to its quotient lattice at $E$.  Taking the
successive root-height pieces gives the two-term determinant
calculation
\[
 \det R\Gamma(X,\mathfrak n^{\mathrm{lv}})
       =\det(k_l)\otimes\det(A_l/h)^{-1}.
\]
Its Euler characteristic is $f_l-e_l$, proving both assertions in
\eqref{eq:T12}.  Finally, the residue calculation of \eqref{eq:W1}
identifies the character on this presentation with $\psi_0$.
\end{proof}

The construction of Lemma~\ref{lem:unipotent-presentation} uses the
root-height filtration and the underlying root bundles.  Consequently
it also applies to the family with bracket $t[-,-]$: the same
$H^1$-vanishing bounds and jet orders work over the auxiliary affine
line.  These geometric presentations and their untwisted structure
maps are equivariant for root height and for the simultaneous scaling
$t\mapsto ct$, $x\mapsto c^{-1}x$ in exponential coordinates.
Under the latter scaling the character has weight $-1$.  The
equivariance just stated will be used for the untwisted degeneration
in Section~\ref{sec:ran}.

\subsection{Localization at a finite quotient}

Choose a sufficiently deep normal congruence subgroup $J\subset L_l$
contained in $H$, with $K_l$ still embedded in $L_l/J$.  Such
subgroups are obtained by truncating the enlarged height coordinates
along the divisor.  For finitely many lattices we use a common
subgroup normal in all of them.  Bars denote the finite quotients;
in particular,
\[
 \bar Z=K_l\backslash\bar L_l,
 \qquad B_N=[\bar Z/\bar H].
\]
For an ordinary smooth module $M$, let $M[J]$ be the submodule killed
by $\operatorname{Lie}J$.  Normality makes it an
$(A_l,H)$-module, and the $M[J]$ exhaust $M$.  Initially, for a bounded
complex, apply this construction termwise.

The finite-action localization of $M[J]$ is the object
$\mathcal V_l(M[J])$ descending from the left-module formula
\begin{equation}\label{eq:T13}
 \mathscr D_{\bar Z}\otimes^L_{U(\bar A_l)}M[J].
\end{equation}
The operator map uses infinitesimal right translations on $\bar Z$;
Harish--Chandra equivariance supplies descent along $\bar H$.
This is the transitive action-algebroid localization of
\eqref{eq:T2}.

If $J'\subset J$, the map $\bar Z'\to\bar Z$ is a torsor under
$J/J'$.  The new action algebroid is the inverse image of the old
one: the added Lie quotient is exactly the relative tangent of this
torsor.  Inflation of the module is therefore pullback as an
action-algebroid module, and \eqref{eq:T13} pulls back to the formula
for $J'$.  For $l\leq l'$, choose a common $J$ and use the inclusion
of the action algebroids on the common presentation of the framed
space.  These maps supply the transition maps in both indices.

Taking the colimit first over $J$ and then over $l$ gives
$\operatorname{Loc}_{N,E}^{\mathrm{lv}}(M)$.  To check the claim,
pull to a chart with a full frame.  Formula~\eqref{eq:T13} there is
the vector/Lie-pair formula for framed changes with Laurent component
in $A_l$.  Replacing a frame modulo $J$ by a full frame only supplies
the components in $\operatorname{Lie}J$.  These components are
uniquely determined because $J$ acts freely, and they act trivially
on $M[J]$.  The stabilizer Lie kernel thus projects isomorphically to
the kernel in the finite action-algebroid calculation.  Increasing
$l$ exhausts the out-current kernel in \eqref{eq:T2}.

This also explains why the construction is valid on bounded derived
inputs.  For a fixed $l$ the kernel on a framed chart has finite rank
$f_l$.  Its Chevalley complex commutes with the exhaustive colimit
of the termwise modules $M[J]$.  The resulting bounded-input
calculation is therefore the ordinary derived Chevalley calculation
with coefficients in $M$.  Exactness of the operation $M\mapsto
M[J]$ at an individual $J$ is not needed.  The comparison to
\eqref{eq:T2} is the colimit of the corresponding maps on out-current
chains.

\subsection{The determinant and cohomological shift}

We now calculate the period of \eqref{eq:T13}.  This is the step
which determines both the normalization of the map and its direction.

\begin{lemma}\label{lem:unipotent-finite-period}
After applying $b$, the finite-action localization satisfies
\begin{equation}\label{eq:T14}
 \begin{aligned}
 &(q_N^{\mathrm{lv}})_*
   \bigl(\operatorname{Exp}(-\psi_0)
                 \otimes^!\mathcal V_l(M[J])\bigr)\\
 &\quad\simeq l^{\mathrm{lv}}[\delta]\otimes
 C^\bullet\!\left(\bar A_l,
    M[J]\otimes\mathcal O(\bar Z)_{-\chi}\right)
       \otimes\det(A_l/h)[e_l].
 \end{aligned}
\end{equation}
Functions carry the derivative action of the right translations;
the subscript $-\chi$ adds the indicated Lie character.  The
isomorphism is natural for congruence refinement and for enlargement
of the height lattice.
\end{lemma}

\begin{proof}
Let $z=\dim(\bar Z/S)$ and $a=\operatorname{rank}\bar A_l$.
Compute ordinary relative left de Rham cohomology on $\bar Z$ and
then shift by $[2\delta]$.  Smooth pullback along the
$\bar H$-presentation and unipotent de Rham acyclicity identify this
with the stated stack period.  Ordinary left de Rham cohomology uses
right tensor with $\Omega^{\mathrm{top}}_{\bar Z}$, with the
exponential twist, and shift $[-z]$.  Since $\bar Z$ is affine,
\eqref{eq:T13} thus gives Lie chains over $U(\bar A_l)$ with this
right coefficient module.

Use the invariant volume line
$\det(\bar A_l/k_l)^{-1}$.  Converting right Lie derivatives on
functions times volume to left derivatives produces the character
$-\chi$.  The groups are unipotent, so their determinant actions
have no modular character.  Finite-dimensional Lie Poincar\'e
duality converts chains to cochains, contributing
$\det(\bar A_l)[a]$.  The resulting line is
\[
 \det(\bar A_l/k_l)^{-1}\otimes\det(\bar A_l)
       =\det(k_l)
       =l^{\mathrm{lv}}\otimes\det(A_l/h).
\]
Because $z=a-f_l$ and $\delta=e_l-f_l$, the resulting shift is
\[
 2\delta-z+a=2\delta+f_l=\delta+e_l.
\]
This proves the line and shift in \eqref{eq:T14}.

It remains to identify the transition maps under this conversion.
For height enlargement the map before Poincar\'e duality is the
inclusion of Lie chains.  Afterwards it is the cochain inclusion
followed by wedging to top degree in the added lattice, with the
relative determinant and shift displayed in \eqref{eq:T14}.  On a
common finite presentation, the determinant comparisons follow from
the exact sequences for the tangent spaces.  Equivalently, the
meromorphic-lattice sequences and the vanishing of the enlarged
root-bundle $H^1$ identify $A_{l'}/A_l$ with $k_{l'}/k_l$ in these
determinant comparisons.  We always use graded determinant lines,
so the wedge maps include the signs of the graded Poincar\'e
conversion.

For congruence refinement the map is inflation of cochains together
with pullback of functions.  One can see this directly from the
Spencer calculation.  If $A$ is a transitive action algebroid on a
smooth scheme and $I=\ker(A\to T)$, its sheaf form is
\begin{equation}\label{eq:unipotent-spencer}
 \operatorname{DR}^{\mathrm{ord}}
   \bigl(\mathscr D\otimes^L_{U(A)}R\bigr)
 \simeq C^\bullet(A,R\otimes\det I)
                           [\operatorname{rank}I].
\end{equation}
Here $\det I$ carries adjoint transport.  Right tensor with top
forms gives Spencer chains; the identity
$\det A\otimes\Omega^{\mathrm{top}}=\det I$ then gives
\eqref{eq:unipotent-spencer}.  The Lie derivative on the two
determinant factors is exactly adjoint transport on $\det I$.
For smooth inverse image of the algebroid, the relative tangent
contributes ordinary smooth de Rham pullback.  Locally in smooth
coordinates this is the product Spencer formula, with ordinary
pullback of forms in the new coordinates.  The Spencer transfer
map is canonical, so these local descriptions glue.  Applying it
to $\bar Z'\to\bar Z$ proves the asserted inflation formula.
The same argument applies to derived coefficient modules: the
bimodule and Spencer resolutions use locally free tangent and Lie
sheaves, and their maps are defined before derived tensor.
\end{proof}

\subsection{The period on bounded inputs and the constant-function map}

After applying $b$, let $M$ be a bounded ordinary Harish--Chandra
complex over the fixed affine parameter scheme $S$.  Take the
double colimit in \eqref{eq:T14}.  For fixed $l$, refinement
of $J$ gives continuous cochains for the compact Lie algebra $A_l$,
using its discrete dual.  Height enlargement gives the semi-infinite
wedge relative to $h$: compact modes contribute cochains, while
modes added outside $h$ contribute the determinant conversion of
chains.  This identifies the bounded ordinary period calculation as
\begin{equation}\label{eq:unipotent-coefficient-period}
 \begin{aligned}
 \operatorname{Per}_{N,E}^{\mathrm{lv}}(M)
 &\simeq l^{\mathrm{lv}}[\delta]\otimes{}\\
 &\quad C^{\infty/2,h}\!\left(
   \operatorname{Lie}LN_E,
   M\otimes\mathcal O(\widetilde B_N)\otimes k_{-\chi}\right).
 \end{aligned}
\end{equation}
Here the functions on the full framed space are the colimit of the
functions on its finite congruence presentations.

This identification includes the differentials.  Choose a
height-compatible topological basis of root modes.  The finite
Chevalley differentials have their usual action and bracket terms.
Congruence inflation is a map of complexes because the congruence
Lie algebras are ideals in $A_l$.  For height enlargement, the terms
outside the old lattice vanish after wedging to top degree: the
old lattice is a subalgebra, and its trace on the intervening quotient
is zero.  Indeed brackets strictly raise root height, so all the
normal-ordering trace terms vanish, including those on lattice
quotients.  The differential in the colimit is therefore the
normal-ordered semi-infinite differential.  For bounded complexes
the ordinary totalization is quasi-isomorphism invariant, since the
coefficient-complex direction is bounded.

The inclusion
\[
 \mathcal O_S\longrightarrow\mathcal O(\widetilde B_N)
\]
of constant functions is equivariant for the loop action.  Applying
semi-infinite cohomology gives \eqref{eq:T10} on bounded ordinary
inputs.  We have thus constructed the map and checked its line,
shift, and congruence transitions.  To obtain the assertion for
families, we must still verify that the construction descends through
the derived presentations and is independent of the formal lift.

\subsection{Derived base change and lift comparisons}

We extend this bounded-input comparison and its constant-function
map to the renormalized category, and then verify its crystalline
compatibilities.
For fixed $l$, use the renormalized Harish--Chandra category of
$(A_l,H)$ as the colimit of its finite-quotient presentations, as in
the compact-Lie calculation in the proof of
Proposition~\ref{loop:strong-action}.  Operations on arbitrary
unbounded inputs will be defined through this categorical colimit
and continuous extension; the termwise $M[J]$ exhaustion above is
used only to identify the bounded ordinary calculation.
On a fixed finite
quotient, construct \eqref{eq:T13} and \eqref{eq:T14} with derived
tensor and the Spencer resolution.  On the cochain side allow all
further congruence refinements, with the pulled function spaces.
Lemma~\ref{lem:unipotent-finite-period} shows that this computes the
same period.  Also form the continuous cochain calculation without
the function space, again as a colimit over further refinements.
The latter transition maps need not be isomorphisms.  Constants give
a natural transformation from this second calculation to the first.

Every operation just described commutes with the relevant derived
base extensions on compact presentations.  At a finite quotient,
the Chevalley and Spencer resolutions have relatively flat terms;
the unipotent-equivariant descent calculations have the same property.
Thus the assertions follow from derived tensor and then colimit.
In the parameter diagrams one uses finite-Tor-dimension changes
between smooth bases, and flat lift changes on infinitesimal
neighborhoods.  This construction does not assert that a fixed
ordinary invariant submodule $M[J]$ commutes with a nonflat base
change.

Compatibility with changing $l$ holds at the same derived level.
For two lattices, compute on a finite quotient of a common larger
lattice and choose a common normal $J$.  Restriction to the smaller
Lie algebra is ordinary derived restriction followed by the
finite-quotient operation.  It commutes with base extension even
when its value is not compact.  For example, restriction of a free
inducing generator has a PBW filtration for the finite-rank difference
of the lattices.  Each term is a module over the smaller Lie algebra,
and their exhaustion computes the restriction and commutes with
base extension.  Equivalently, this is restriction for the induction
monads.  The algebroid inclusion and the top-wedge maps of
Lemma~\ref{lem:unipotent-finite-period} are defined on this common
finite quotient and therefore remain compatible after colimit.

Restriction from the original loop Harish--Chandra category to
$(A_l,H)$ is likewise restriction of induction-monad modules.  The
equivariance category is the same on both sides, so the monadic
calculation of Proposition~\ref{loop:hc-renormalization} makes this
restriction continuous and $t$-exact after applying $b$.  Its agreement with the bounded construction can be checked
on an inducing generator.  If $W$ is its locally free inducing
representation, then
\[
 \operatorname{Ind}_{h}^{\operatorname{Lie}LN_E}W
   =\mathop{\operatorname{colim}}_{l'\geq l}
                         \operatorname{Ind}_{h}^{A_{l'}}W.
\]
For a fixed $l'$, take $J$ normal also in $L_{l'}$ and sufficiently
deep for $W$.  The resulting PBW-flat module factors through the
finite Lie quotient.  It yields precisely the colimit and
constant-function map constructed above.

We now identify the natural transformations in these comparisons.
Each finite presentation maps naturally to the bounded
calculation by inflation and ordinary comparison; for localization
this is the action-algebroid map after pullback to full-frame charts.
On bounded inputs in the loop category, comparison with the
continuous calculation is obtained by evaluating these maps on
maps from compact objects, using the fully faithful bounded
comparison of
Proposition~\ref{loop:hc-renormalization}.  The displayed exhaustion proves that
these comparison maps are isomorphisms on inducing generators,
hence on their thick closure.  Finally, the comparisons with
\eqref{eq:T2} and \eqref{eq:L6} are Lie-algebroid inclusions and
the same wedge identifications.  Their terms on the inducing
generators are locally free or filtered unions of flat modules.
They therefore commute with the derived base changes in question.

It remains to compare lifts.  The full-frame presentation
\eqref{eq:T11} uses the absolute formal curves and is independent of
the lift.  For two lifts, compare cofinal height lattices: a pole
bound coming from either lift is contained in a sufficiently large
bound coming from the other.  Relative divisor thickenings give
smooth lattice differences, including in the nilpotent comparison
diagrams.  With these common enlargements, the presentation by
$K_l$ and the congruence refinements apply on both sides.  All
comparison maps are the inflation and top-wedge maps already proved
compatible in Lemma~\ref{lem:unipotent-finite-period}.

For a change between nonzero levels, choose $J$ sufficiently deep
for both.  The same $\bar Z$ and $\bar A_l$ then calculate the two
periods; only the quotient by the level group changes.  Smooth
pullback, with normalization $[-2\delta]$, identifies their ordinary
de Rham calculations on $\bar Z$.  On the semi-infinite side,
\eqref{eq:T9} is exactly the change of lattice.  The determinant and
dimension formulas \eqref{eq:T12} identify the two normalizations.
These comparisons are transitive, including when a common deeper
level is needed.  We obtain a map of crystalline operations.
Continuous extension from compacts completes the proof of
Proposition~\ref{prop:unipotent-period}.

\subsection{Restriction from the reductive group}

To apply the period map to a $G$ input, one must compare the
continuous $N$ calculation after restriction with the bounded BRST
calculation used to define the $G$ operation.  This comparison is
needed because restriction need not preserve compact objects.

\begin{lemma}\label{lem:unipotent-restriction}
With the restriction functor of Section~\ref{sec:localization}, there
is a natural isomorphism of continuous operations
\begin{equation}\label{eq:T15}
 B_E^0\circ\operatorname{Res}\simeq B_E^0,
\end{equation}
where the right side denotes the operation on the $G$ input category.
Consequently \eqref{eq:T10} also gives the period map for $G$ inputs,
with the same parameter and level compatibilities.
\end{lemma}

\begin{proof}
On compact $G$ inputs, restriction is computed by bounded ordinary
modules.  On a bounded $N$ module there is a natural comparison from
the continuously extended $B_E^0$ to the bounded BRST formula:
evaluate the latter on maps from compacts and use the fully faithful
bounded comparison of
Proposition~\ref{loop:hc-renormalization}.  This map is the identity on
compact $N$ inputs.

For restriction of an inducing $G$ generator, use the outside-mode
PBW filtration \eqref{eq:T4}.  Each graded term is induced from a
representation exhausted in heart degrees by finite representations.
Bounded BRST complexes commute termwise with these filtered colimits
in a bounded range of coefficient degrees.  For derived filtered
diagrams in the same range, finite truncation triangles reduce this
assertion to the assertion in the heart.  The comparison therefore
computes the exhausted graded terms and then the filtration itself.
It is an isomorphism on restricted inducing generators and hence on
the thick category of compact $G$ inputs.  This proves
\eqref{eq:T15}, including naturality, by continuous extension.

Composing \eqref{eq:T10} with the localization comparison
\eqref{eq:T5} now gives the asserted $G$ map.  The preceding proof
also justifies the bounded BRST calculations on inserted exhausted
modules used below.  It makes no comparison of ordinary and
renormalized formulas on arbitrary unbounded complexes.
\end{proof}
